% Standalone build wrapper for data-management-plan.tex.
% Produces data-management-plan.pdf directly via pdflatex (never via Preview/Quartz
% export, which corrupts PDF object dictionaries and gets rejected by research.gov).
% Build with: make data-management-plan
%
% NSF 25-515 (SaTC 2.0) Data Management and Sharing Plan rules: Supplementary
% Document, no more than two pages, must be substantive and project-specific.
% Beyond standard PAPPG content, SaTC 2.0 asks proposers to address: (1) ethics
% of data collection and harms that could arise from sensitive data; (2) methods
% for controlled-access data sharing; (3) security techniques preventing
% unauthorized dissemination; and (4) whether data is LLM/synthetic or produced
% by real experiments, tagged for provenance, accountability, and reproducibility.
\documentclass[11pt,onecolumn]{article}
\usepackage[letterpaper,margin=1in]{geometry}
\usepackage{times}
\usepackage[colorlinks=true,linkcolor=blue,citecolor=blue,urlcolor=blue]{hyperref}
\usepackage{xcolor}
\usepackage{soul}
\linespread{1.0}
\setlength{\parindent}{0em}
\setlength{\parskip}{0.4em}
\pagestyle{empty}

\newcommand{\arman}[1]{%
	\begingroup
	\definecolor{hlcolor}{RGB}{135, 216, 230}\sethlcolor{hlcolor}%
	\textcolor{black}{\hl{\textit{\textbf{Arman:} #1}}}%
	\endgroup
}

\newcommand{\head}[1]{\noindent{\bf #1}}

\begin{document}

\section*{Data Management and Sharing Plan}

This plan describes how the data, software, and educational materials produced by
VEST: Verifiable Trust Establishment on Embedded Systems will be generated,
documented, shared, and preserved, consistent with NSF's public access and
data-sharing policy.

\head{Data and Software.}
The project produces four kinds of output, one per research thrust plus a shared
evaluation layer.
\emph{Thrust 1 (trusted-domain composition and attestation):} FPGA enclave
configuration artifacts and measurement logs from the BYOTee toolchain; the
Root-monitor firmware and domain-manifest tooling developed for ARM's Confidential
Compute Architecture on the Fixed Virtual Platform; and the resulting hardware and
software measurement records used to evaluate composition verifiability and
trusted-computing-base size.
\emph{Thrust 2 (runtime trust attestation):} the LLVM compiler passes and runtime
instrumentation extending the ENOLA control-flow attestation engine to
interrupt-driven and multitasking execution, and the resulting evidence-size,
verifier-latency, and overhead measurements collected on the ARM Versatile
Express V2M-MPS3 board.
\emph{Thrust 3 (compartmentalization and recovery):} the compiler-driven
partitioning tool, transition-monitor implementation, and data-flow tagging
runtime, and the resulting MPU-switching overhead, violation-detection-latency,
and containment/recovery-time measurements collected on the V2M-MPS3 and
RA8M1 (EK-RA8M1) boards.
\emph{Shared evaluation data:} benchmark results and comparisons against
established datasets (BEEBS, BenchIoT) and against the prior work discussed in
each thrust's Related Work.
\emph{Educational materials:} lab modules and challenge material developed for
the PI's undergraduate and graduate systems-security course and for the MITRE
Embedded Capture-the-Flag team.
All research data in this project is generated directly by hardware and software
experiments described in the Evaluation Plan; none of it is synthetic or
produced by a large language model, and every record is tagged with the
generating platform (FVP, V2M-MPS3, or EK-RA8M1), build configuration, and
software commit, so its provenance is traceable without additional
documentation.

\head{Standards and Formats.}
Source code and compiler passes are written in C and LLVM IR; firmware and
monitor code targets the TF-A/TF-RMM and Cortex-M toolchains used in the
Evaluation Plan. Measurement data (evidence size, latency, overhead, and
trusted-computing-base size) is stored as CSV with a documented schema
describing each field and the experiment that produced it. Configuration
artifacts (FVP platform configs, board BSPs, MPU/GPT manifests) are stored as
plain text or JSON alongside the code that consumes them, so an experiment can
be reconstructed without guessing undocumented state. Documentation is
Markdown and PDF.

\head{Access, Sharing, and Provenance.}
All data, software, and educational materials produced by this project are
released without access restriction: none of VEST's outputs are attack
tooling, personally identifiable, or export-controlled, so nothing in this
project requires controlled or restricted-access sharing.
GitHub is the working repository for code, data, and documentation; tagged
releases are mirrored to Zenodo with DOIs for long-term citability, and
publications are deposited through NSF's public access repository.
\arman{Confirm the GitHub organization/repository name(s) and Zenodo community
link before submission.}
The one circumstance under which release would be delayed is if evaluation on
a third-party reference implementation (for example, ARM's Trusted Firmware-A
or TF-RMM) incidentally surfaces a genuine vulnerability; in that case the
finding is held privately and disclosed to the maintainer under standard
coordinated-disclosure practice before publication, while the rest of the
project's outputs are released on schedule.
Software repositories and datasets are protected against unauthorized
modification using GitHub's access controls and branch protection, with
releases signed and tagged so a downstream user can verify provenance.

\head{Reuse and Redistribution.}
Software is released under a permissive open-source license (MIT or BSD), and
data and educational material under a Creative Commons license, so other
researchers, educators, and practitioners can reuse, redistribute, and build
on them with attribution.
Where project code extends third-party open-source components, such as ARM's
Trusted Firmware-A and TF-RMM reference implementations or the LLVM toolchain,
redistribution respects that project's own license; only the instrumentation,
monitor, and analysis code developed in this project is newly licensed.
Version control and a public issue tracker on GitHub support responsible reuse
and let other groups extend or reproduce the project's results independently
of the PI's lab, consistent with the reproducibility goals stated throughout
the proposal.

\head{Archiving and Preservation.}
Data, code, and documentation are retained for at least three years beyond the
end of the award on backed-up storage at UTEP under institutional IT policy,
and mirrored indefinitely on GitHub and Zenodo so long-term access does not
depend on the university's retention window alone. Publications remain
accessible through the publisher's repository and NSF's public access
repository.

\head{Ethics and Human Subjects.}
This project involves no human subjects and no personally identifiable
information: every dataset is a hardware or software measurement collected
from equipment the PI's lab owns or controls, or from public reference
implementations evaluated in emulation, so IRB review is not required. Student
researchers who contribute to the project do so as project personnel and,
where appropriate, co-authors, not as human-subject participants.

\head{Prior Experience.}
The PI has a record of releasing research artifacts responsibly: both BYOTee,
the FPGA-based trusted-execution framework Thrust~1 builds on, and ENOLA, the
control-flow attestation system Thrust~2 builds on, have been open-sourced
with their implementations and evaluations publicly released. The PI's
capture-the-flag and course material has likewise been shared with the
community he teaches and competes in.

\end{document}
